% GENERATED FILE. Edit canonical lessons, then run npm run generate:books.
% canonical-source-hash: fnv1a64:18cba365566a4f2a
% canonical-lessons: RU-C304-detstvo, RU-C304-yunost, RU-C304-molodost, RU-C304-starost, RU-C304-vzroslet

\chapter{Stages of life}
\label{ch:ru-stages-of-life}
\hlchaptermodality{304}

\begin{chapteropening}
\textbf{By the end of this chapter:} \emph{I can talk about childhood, the teenage years, youth and old age, and say that someone grew up.}

Complete the last lesson of chapter 304: I can talk about childhood, the teenage years, youth and old age, and say that someone grew up.
\end{chapteropening}

\section[\texorpdfstring{\ru{детство} (détstvo)}{детство (détstvo)}]{\ru{детство} (détstvo) — childhood}

\label{lesson:RU-C304-detstvo}



\begin{quote}
Before the new one: say the Russian for pregnant, then the Russian for to give birth.
\end{quote}

\subsection*{You'll want to know: \ru{детство}}
\textbf{\ru{детство}} (\emph{détstvo}) — \textquotedblleft{}childhood\textquotedblright{}.

\begin{quote}
\ru{Детство} \ru{я} \ru{провёл} \ru{в} \ru{деревне}.
\end{quote}

\begin{quote}
\emph{Détstvo ya provyól v derévne.}
\end{quote}

\begin{quote}
I spent my childhood in the countryside.
\end{quote}

Neuter. \emph{\ru{С} \ru{детства}} — since childhood.

\begin{cousinweb}
\emph{\ru{детство}} is built on \emph{\ru{дети}} \textquotedblleft{}children\textquotedblright{}, from Proto-Slavic \emph{dětę} \textquotedblleft{}a child\textquotedblright{}, which EtymDB traces to PIE \emph{d\textsuperscript{h}eh\textsubscript{1}(y)-} \textquotedblleft{}to suckle\textquotedblright{}.
\end{cousinweb}

\begin{grammarlens}[title={What you've built}]
One more word for telling, explaining and describing at B1.
\end{grammarlens}

\subsection*{Your turn}
\begin{itemize}
\raggedright
  \item \emph{Say it:} \emph{détstvo}
  \item \emph{Say it:} \emph{détstvo}, once more
  \item \emph{Say it:} a sentence of your own with \emph{détstvo}
  \item \emph{Recall:} say the Russian for pregnant, then the Russian for to give birth, then say \emph{détstvo} again
\end{itemize}

\begin{tcolorbox}[breakable,colback=teal!4,colframe=teal!35!black,title={Before you move on}]
What does \ru{детство} mean? (\textquotedblleft{}Childhood\textquotedblright{}.) Say it once more.
\end{tcolorbox}
\section[\texorpdfstring{\ru{юность} (yúnost)}{юность (yúnost)}]{\ru{юность} (yúnost) — youth, the teenage years}

\label{lesson:RU-C304-yunost}



\begin{quote}
Before the new one: say the Russian for to give birth, then the Russian for childhood.
\end{quote}

\subsection*{You'll want to know: \ru{юность}}
\textbf{\ru{юность}} (\emph{yúnost}) — \textquotedblleft{}youth, the teenage years\textquotedblright{}.

\begin{quote}
\ru{В} \ru{юности} \ru{он} \ru{играл} \ru{на} \ru{гитаре}.
\end{quote}

\begin{quote}
\emph{V yúnosti on igrál na gitáre.}
\end{quote}

\begin{quote}
In his youth he played the guitar.
\end{quote}

Feminine. Roughly the years from fifteen to twenty.

\begin{cousinweb}
\emph{\ru{юность}} is built on \emph{\ru{юный}} \textquotedblleft{}young\textquotedblright{}, which EtymDB marks as a borrowing from Church Slavonic \emph{junŭ}, from Proto-Slavic \emph{junŭ} and PIE \emph{h\textsubscript{2}yuh\textsubscript{1}en-} \textquotedblleft{}young\textquotedblright{}; English \emph{young} is a cousin.
\end{cousinweb}

\begin{grammarlens}[title={What you've built}]
One more word for telling, explaining and describing at B1.
\end{grammarlens}

\subsection*{Your turn}
\begin{itemize}
\raggedright
  \item \emph{Say it:} \emph{yúnost}
  \item \emph{Say it:} \emph{yúnost}, once more
  \item \emph{Say it:} a sentence of your own with \emph{yúnost}
  \item \emph{Recall:} say the Russian for to give birth, then the Russian for childhood, then say \emph{yúnost} again
\end{itemize}

\begin{tcolorbox}[breakable,colback=teal!4,colframe=teal!35!black,title={Before you move on}]
What does \ru{юность} mean? (\textquotedblleft{}Youth, the teenage years\textquotedblright{}.) Say it once more.
\end{tcolorbox}
\section[\texorpdfstring{\ru{молодость} (mólodost)}{молодость (mólodost)}]{\ru{молодость} (mólodost) — one's young years, youth}

\label{lesson:RU-C304-molodost}



\begin{quote}
Before the new one: say the Russian for childhood, then the Russian for youth.
\end{quote}

\subsection*{You'll want to know: \ru{молодость}}
\textbf{\ru{молодость}} (\emph{mólodost}) — \textquotedblleft{}one's young years, youth\textquotedblright{}.

\begin{quote}
\ru{В} \ru{молодости} \ru{бабушка} \ru{много} \ru{путешествовала}.
\end{quote}

\begin{quote}
\emph{V mólodosti bábushka mnógo puteshéstvovala.}
\end{quote}

\begin{quote}
In her youth Grandma travelled a lot.
\end{quote}

Feminine. \emph{\ru{Молодость}} reaches into the twenties and thirties; \emph{\ru{юность}} is the teens.

\begin{cousinweb}
\emph{\ru{молодость}} is built on \emph{\ru{молодой}} \textquotedblleft{}young\textquotedblright{} and the ending \emph{-\ru{ость}} (EtymDB); \emph{\ru{молодой}} is inherited from Proto-Slavic \emph{moldŭ} \textquotedblleft{}young\textquotedblright{}.
\end{cousinweb}

\begin{grammarlens}[title={What you've built}]
One more word for telling, explaining and describing at B1.
\end{grammarlens}

\subsection*{Your turn}
\begin{itemize}
\raggedright
  \item \emph{Say it:} \emph{mólodost}
  \item \emph{Say it:} \emph{mólodost}, once more
  \item \emph{Say it:} a sentence of your own with \emph{mólodost}
  \item \emph{Recall:} say the Russian for childhood, then the Russian for youth, then say \emph{mólodost} again
\end{itemize}

\begin{tcolorbox}[breakable,colback=teal!4,colframe=teal!35!black,title={Before you move on}]
What does \ru{молодость} mean? (\textquotedblleft{}One's young years, youth\textquotedblright{}.) Say it once more.
\end{tcolorbox}
\section[\texorpdfstring{\ru{старость} (stárost)}{старость (stárost)}]{\ru{старость} (stárost) — old age}

\label{lesson:RU-C304-starost}



\begin{quote}
Before the new one: say the Russian for youth, then the Russian for one's young years.
\end{quote}

\subsection*{You'll want to know: \ru{старость}}
\textbf{\ru{старость}} (\emph{stárost}) — \textquotedblleft{}old age\textquotedblright{}.

\begin{quote}
\ru{В} \ru{старости} \ru{он} \ru{жил} \ru{у} \ru{дочери}.
\end{quote}

\begin{quote}
\emph{V stárosti on zhil u dócheri.}
\end{quote}

\begin{quote}
In his old age he lived with his daughter.
\end{quote}

Feminine.

\begin{cousinweb}
\emph{\ru{старость}} is inherited from Proto-Slavic \emph{starostĭ}, built on \emph{starŭ} \textquotedblleft{}old\textquotedblright{}, which EtymDB traces to PIE \emph{steh\textsubscript{2}-ro-} \textquotedblleft{}large, old\textquotedblright{}.
\end{cousinweb}

\begin{grammarlens}[title={What you've built}]
One more word for telling, explaining and describing at B1.
\end{grammarlens}

\subsection*{Your turn}
\begin{itemize}
\raggedright
  \item \emph{Say it:} \emph{stárost}
  \item \emph{Say it:} \emph{stárost}, once more
  \item \emph{Say it:} a sentence of your own with \emph{stárost}
  \item \emph{Recall:} say the Russian for youth, then the Russian for one's young years, then say \emph{stárost} again
\end{itemize}

\begin{tcolorbox}[breakable,colback=teal!4,colframe=teal!35!black,title={Before you move on}]
What does \ru{старость} mean? (\textquotedblleft{}Old age\textquotedblright{}.) Say it once more.
\end{tcolorbox}
\section[\texorpdfstring{\ru{взрослеть} (vzroslét)}{взрослеть (vzroslét)}]{\ru{взрослеть} (vzroslét) — to grow up}

\label{lesson:RU-C304-vzroslet}



\begin{quote}
Before the new one: say the Russian for one's young years, then the Russian for old age.
\end{quote}

\subsection*{You'll want to know: \ru{взрослеть}}
\textbf{\ru{взрослеть}} (\emph{vzroslét}) — \textquotedblleft{}to grow up\textquotedblright{}.

\begin{quote}
\ru{Дети} \ru{быстро} \ru{взрослеют}.
\end{quote}

\begin{quote}
\emph{Déti býstro vzrosléyut.}
\end{quote}

\begin{quote}
Children grow up fast.
\end{quote}

Imperfective; perfective \emph{\ru{повзрослеть}} (\emph{povzroslét}). \emph{\ru{Взрослый}} — an adult.

\begin{cousinweb}
\emph{\ru{взрослеть}} is built on \emph{\ru{взрослый}} \textquotedblleft{}adult\textquotedblright{}, from \emph{\ru{вз}-} \textquotedblleft{}up\textquotedblright{} and the root \emph{\ru{рос}-} of \emph{\ru{расти}} \textquotedblleft{}to grow\textquotedblright{}: grown up. EtymDB lists no parent.
\end{cousinweb}

\begin{grammarlens}[title={What you've built}]
One more word for telling, explaining and describing at B1.
\end{grammarlens}

\subsection*{Your turn}
\begin{itemize}
\raggedright
  \item \emph{Say it:} \emph{vzroslét}
  \item \emph{Say it:} \emph{vzroslét}, once more
  \item \emph{Say it:} a sentence of your own with \emph{vzroslét}
  \item \emph{Recall:} say the Russian for one's young years, then the Russian for old age, then say \emph{vzroslét} again
\end{itemize}

\begin{tcolorbox}[breakable,colback=teal!4,colframe=teal!35!black,title={Before you move on}]
What does \ru{взрослеть} mean? (\textquotedblleft{}To grow up\textquotedblright{}.) Say it once more.
\end{tcolorbox}
